% id: 84-15
% book: 베이직쎈 미적분1 · 05 도함수의 활용 (2)
% section: 개념 37 다항함수가 극값을 가질 조건
% passage: 다음 함수 $f(x)$가 극댓값과 극솟값을 모두 갖도록 하는 실수 $a$의 값의 범위를 구하시오.
% figure: none
% answer: 
% answer_source: 
% variant_level: 0 (원본 전사)
% 이 파일은 build.mjs 가 items.json 에서 만든다. 고칠 때는 items.json 을 고친다.
\smash{\raisebox{1.5mm}{\dmkichul{베이직쎈 미적분1 84쪽 15번}}}
$f(x)=x^4-2x^3+ax^2-5$
\dmhint{$f(x)=x^4-2x^3+ax^2-5$에서\\ \hspace*{1.5em}$f'(x)=4x^3-6x^2+2ax=2x(2x^2-3x+a)$\\ 사차함수 $f(x)$가 극댓값과 극솟값을 모두 가지려면 삼차방정식 $f'(x)=0$이 \blank[6em]을 가져야 하므로 이차방정식 $2x^2-3x+a=0$이 \underline{0이 아닌} \blank[6em]을 가져야 한다.\\ \hspace*{1.5em}$\therefore a\ne 0$ \hfill $\cdots\cdots$ ㉠\\ \hspace*{1.5em}{\footnotesize $2x(2x^2-3x+a)=0$에서 $x=0$ 또는 $2x^2-3x+a=0$}\\ 이차방정식 $2x^2-3x+a=0$의 판별식을 $D$라 하면\\ \hspace*{1.5em}$D=9-8a\blank 0$ \qquad $\therefore a\blank\dfrac{9}{8}$ \hfill $\cdots\cdots$ ㉡\\ ㉠, ㉡에서 \blank[7em]}
